\section{Tempo e spazio}

\subsection{Complessitá di tempo}
$M(x)$ richiede tempo $t$ sse:
\[
\configurazionestart{x} \xrightarrow{M^t} \configurazione{H}{w}{u}
\]
Dove $H \in \{h, "yes", "no"\}. $\\

\textbf{$M$ opera in tempo $f(n)$ sse $\forall x \in \Sigma^*$ $M(x)$ richiede tempo $t <= f(n)$} \\

\subsection{Complessitá di spazio}
Si potrebbe considerare la somma della lunghezza dei $k$ nastri: $\sum_{i=2}^{k-1} |w_i u_i|$ \\
Ma possiamo scegliere anche la lunghezza del nastro piú lungo: $\max_{i \in \{1, \dots, k\}} |w_i u_i|$. \\
Queste due quantitá sono asintoticamente equivalenti:\\
$\max_{i \in \{1, \dots, k\}} |w_i u_i| \leq  \sum_{i=1}^{k} |w_i u_i| \leq k \cdot \max_{i \in \{1, \dots, k\}} |w_i u_i|$  \\

\textbf{Attenzione:} con queste definizioni naive, lo spazio occupato dall'input verrá sempre contato.
Questo vuol dire che lo spazio é sempre limitato inferiormente da $n$, anche quando non usiamo variabili ausiliarie. \\
A noi interessa considerare lo \textit{spazio ausiliario} utilizzato. Per farlo, occorre considerare le macchine di Turing con input e output e contare solo i nastri interni. \\

\subsubsection{Complessitá di spazio ausiliario nelle MdT con input e output}
$M(x)$ richiede spazio $z$ sse: \\
Se:
\[
\configurazionestart{x} \xrightarrow{M^*} \configurazione{H}{w}{u}
\]
Dove $H \in \{h, "yes", "no"\}$, allora: \\
$z = \sum_{i=2}^{k-1} |w_i u_i|$ \\

\textbf{$M$ opera in spazio $g(n)$ sse $\forall x \in \Sigma^*$ $M(x)$ richiede spazio $z <= g(n)$} \\
